\section{Results and Discussion}
\label{sec:results}

Experimental results are not yet available. No performance numbers are reported here. The following subsections will be completed after the planned experiments.

\subsection{Quantitative Results}
[TO BE COMPLETED AFTER EXPERIMENTS]

Table~\ref{tab:comparison} will summarize baseline and proposed-model performance, emphasizing sensitivity and specificity in addition to accuracy, F1-score, and ROC-AUC.

\begin{table}[!t]
\caption{Comparison of baseline and neuro-symbolic models (placeholder).}
\label{tab:comparison}
\centering
\small
\begin{tabular}{@{}lcccc@{}}
\toprule
Model & Acc. & Sens. & Spec. & F1 \\
\midrule
Neural baseline & --- & --- & --- & --- \\
Concept neural baseline & --- & --- & --- & --- \\
Proposed (DeepProbLog) & --- & --- & --- & --- \\
\bottomrule
\end{tabular}
\end{table}

\subsection{Comparison with Neural Baselines}
[TO BE COMPLETED AFTER EXPERIMENTS]

This subsection will discuss whether the neuro-symbolic model improves, matches, or underperforms the neural baselines. Accuracy improvement remains a research hypothesis.

\subsection{Interpretability and Error Analysis}
[TO BE COMPLETED AFTER EXPERIMENTS]

Planned analyses include inspection of reasoning traces for selected cases, assessment of concept-prediction failures that propagate into symbolic inference, and identification of settings where medical knowledge constraints help or hurt the final decision. Availability of an explicit trace will be distinguished from clinical adequacy of that trace.
